\section{Additional Exercises}

\subsection{Pull-Back \& Push-Forward}

\begin{exercise}\label{ex:symmetry-1}
Consider a smooth map $\phi:\cM\to\cN$ between two differential manifolds. Show that a function $f\in C^{\infty}(\cN)$, the pull-back of the gradient of $f$ is the same as the gradient of the pull-back of $f$, i.e.
\bse
    \phi^*(df) = d(\phi^*f).
\ese
\end{exercise}
\hyperref[sol:symmetry-1]{\small \textit{Solution on p.~\pageref*{sol:symmetry-1}}}

\begin{exercise}\label{ex:symmetry-2}
The push-forward $\phi_* : T\cM \to T\cN$ is a linear map between tangent bundles. Calculate its component functions
\bse
    \phi_{*\,\, b}^{\,\, a} := dy^a : \phi_*\bigg(\frac{\p}{\p x^b}\bigg)
\ese
with respect to charts $(U\ss\cM, x)$ and $(V\ss \cN, y)$!
\end{exercise}
\hyperref[sol:symmetry-2]{\small \textit{Solution on p.~\pageref*{sol:symmetry-2}}}

\begin{exercise}\label{ex:symmetry-3}
Show that the component functions of the pull back $\phi^*g$ of the metric tensor field are obtained from the component functions of $g$ by
\bse
    (\phi^*g)_{ab}(p) = \bigg(\frac{\p(y\circ\phi)^m}{\p x^a}\bigg)_p \bigg(\frac{\p(y\circ\phi)^n}{\p x^b}\bigg)_p g_{mn}\big(\phi(p)\big).
\ese
\end{exercise}
\hyperref[sol:symmetry-3]{\small \textit{Solution on p.~\pageref*{sol:symmetry-3}}}

\subsection{Lie Derivative --- The Pedestrian Way}

\begin{exercise}\label{ex:symmetry-4}
Consider the smooth embedding $\iota:S^2\to \R^3$ of $(S^2,\cO,\cA)$ into $(\R^3,\cO_{st},\cB)$, which for the familiar chart $(U,x)\in\cA$ and $(\R^3,y=\b1_{\R^3})\in\cB$ is given by
\bse
    y\circ\iota\circ x^{-1} : (\vartheta,\varphi) \mapsto (a\cos\varphi\sin\vartheta, b\sin\varphi\sin\vartheta,c\cos\vartheta),
\ese
where $a,b$ and $c$ are positive real numbers. What can you say about the shape of $\iota(S^2)$?
\end{exercise}
\hyperref[sol:symmetry-4]{\small \textit{Solution on p.~\pageref*{sol:symmetry-4}}}

\begin{exercise}\label{ex:symmetry-5}
Now assume $(\R^3,\cO_{st},\cB)$ is additionally equipped with the Euclidean metric $g$, whose components with respect to the chart $(\R^3,y)$ are given by
\bse
    g_{ab}(p) = \begin{pmatrix}
    1 & 0 & 0 \\
    0 & 1 & 0 \\
    0 & 0 & 1
    \end{pmatrix} \qquad \text{for any } p\in U.
\ese
Write down the component functions of $g^{\text{ellipsoid}}:=\iota^*g$ with respect to the chart $(U,x)$!
\end{exercise}
\hyperref[sol:symmetry-5]{\small \textit{Solution on p.~\pageref*{sol:symmetry-5}}}

\begin{exercise}\label{ex:symmetry-6}
For convenience, denote by $(\vartheta,\varphi)$ the coordinate functions $(x^1,x^2)$. Check the vector fields
\bse
    \begin{split}
        X_1(p) & = -\sin\varphi(p)\bigg(\frac{\p}{\p \vartheta}\bigg)_p - \cot\vartheta(p)\cos\varphi(p)\bigg(\frac{\p}{\p \varphi}\bigg)_p \\
        X_2(p) & = \cos\varphi(p)\bigg(\frac{\p}{\p \vartheta}\bigg)_p - \cot\vartheta(p)\sin\varphi(p)\bigg(\frac{\p}{\p \varphi}\bigg)_p \\
        X_3(p) & = \bigg(\frac{\p}{\p\varphi}\bigg)_p
    \end{split}
\ese
constitute a Lie subalgebra of $(\Gamma TS^2,[\cdot,\cdot])$ and determine the structure constants!
\end{exercise}
\hyperref[sol:symmetry-6]{\small \textit{Solution on p.~\pageref*{sol:symmetry-6}}}

\begin{exercise}\label{ex:symmetry-7}
Calculate the integral curve of $X_3$ through the point $p=x^{-1}(\vartheta_0,\varphi_0)$, i.e. the curve $\gamma_p$ satisfying
\bse
    \gamma_p(0) = p, \qand v_{\gamma_p,\gamma_p(\lambda)} = (X_3)_{\gamma_p(\lambda)}
\ese
in the chart $(U,x)$!
\end{exercise}
\hyperref[sol:symmetry-7]{\small \textit{Solution on p.~\pageref*{sol:symmetry-7}}}

\begin{exercise}\label{ex:symmetry-8}
The integral curves $\gamma_p$ give rise to a one-parameter family of smooth maps $h^{X_3}_{\lambda}:S^2\to S^2$. Calculate the pull-back
\bse
    \Big(h^{X_3}_{\lambda}\Big)^*g^{\text{ellipsoid}}
\ese
of the metric on $S^2$. What can you conclude about the Lie derivative $\cL_{X_3}g^{\text{ellipsoid}}$?
\end{exercise}
\hyperref[sol:symmetry-8]{\small \textit{Solution on p.~\pageref*{sol:symmetry-8}}}
